\documentclass[11pt,a4paper]{article}
% Shared preamble for RuPhO English problem statements (match T1 2026 style).
% Use: \input{latex/rupho_en_preamble.tex} after \documentclass.
\usepackage[margin=1in]{geometry}
\usepackage{amsmath,amssymb}
\usepackage{graphicx}
\usepackage{float}
\usepackage{enumitem}
\usepackage{microtype}
\usepackage{caption}
\usepackage{tabularx}
\usepackage{hyperref}

\captionsetup{font=small,labelfont=bf,skip=6pt}

\setlist[itemize]{leftmargin=1.2cm}
\setlist[enumerate]{leftmargin=1.2cm}

% One outer box per subproblem: [id | statement | points] (no inner rules)
\newcommand{\problembox}[3]{%
  \par\addvspace{0.42em}%
  \noindent
  \setlength{\fboxsep}{7.5pt}%
  \setlength{\fboxrule}{0.95pt}%
  \fbox{%
    \begin{minipage}{\dimexpr\linewidth-2\fboxsep-2\fboxrule\relax}
    \begin{tabularx}{\linewidth}{@{}>{\raggedright\arraybackslash\bfseries}p{2.1em} @{\hspace{0.9em}} >{\raggedright\arraybackslash}X @{\hspace{0.9em}} >{\raggedleft\arraybackslash\bfseries}p{2.4em}@{}}
    #1 & #2 & #3
    \end{tabularx}
    \end{minipage}%
  }%
  \par\addvspace{0.32em}%
}


\title{\textbf{Irreversible thermodynamic cycles}}
\author{\normalsize Y25 (2025) --- English translation}
\date{}
\begin{document}
\maketitle

It is well known that the maximum efficiency of a heat engine operating with two thermal reservoirs (heater and refrigerator), at fixed temperatures of the heater and refrigerator, corresponds to the Carnot cycle. However, in this cycle, heat transfer from the heater to the working fluid occurs at equal temperatures, and therefore infinitely slowly. In any real engine, a finite temperature difference must be maintained during heat transfer, which leads to a decrease in efficiency. In this task we will consider different aspects of the behavior of engines, taking into account the fact that processes do not occur infinitely slowly, which means that processes in the engine are irreversible.

\section*{Part A. Engine with friction (1.5 points)}

In this part, we will assume that in addition to performing useful work, the engine piston overcomes the force of friction. We will assume that the engine piston moves in a cylinder of constant cross-section along the guides of this cylinder. Suppose that during a cycle the piston moves a distance $D$ and back, a quantity of heat $Q_1>0$ is supplied to it, and a quantity of heat $Q_2 >0$ is removed. The friction force depends on the piston speed as $F = \gamma v$, where $v$ is the piston speed. Cycle duration $\tau$. For simplicity, we will assume that for half the cycle the piston moves at a constant speed in one direction, and for the second half - in the other direction.

\problembox{A1}{Write down an expression for the net power of the engine. Express your answer in terms of $Q_1$, $Q_2$, $\gamma$, $D$, $\tau$.}{0.70}

\problembox{A2}{Find the optimal cycle duration $\tau$ at which the power will be maximum.}{0.30}

\problembox{A3}{Find an expression for efficiency (the ratio of useful work to the amount of heat supplied) when the engine operates at the optimal cycle time.}{0.30}

\problembox{A4}{Let $Q_1$, $Q_2$ be the same as if the engine were operating according to an ideal Carnot cycle. Express the efficiency of the cycle under consideration in terms of the efficiency of such a Carnot cycle $\eta_c$.}{0.20}

\section*{Part B. Engine taking into account finite thermal conductivity (5.5 points)}

Let the temperatures of the heater and cooler of the heat engine be $T_1$, $T_2$, respectively. The working fluid is $\nu$ moles of an ideal monatomic gas. By analogy with the Carnot cycle, the process consists of two isotherms with temperatures $T_1'$ and $T_2'$, with $T_1 > T_1' > T_2' > T_2$ connected by two adiabats. Throughout Part B, temperatures $T_1$ and $T_2$ are considered constant. We will consider the heat transfer power to be proportional to the temperature difference $$ P_1 = \textbackslash\{\}alpha_1 \textbackslash\{\}Delta T _1 = \textbackslash\{\}alpha_1 (T_1 - T_1'), \textbackslash\{\}quad P_2 = \textbackslash\{\}alpha_2 \textbackslash\{\}Delta T_2= \textbackslash\{\}alpha_2 (T_2' - T_2). $$На изотерме с температурой $T_1'$ к газу подводится количество тепла $Q_1$, а его объем увеличивается от $V_1$ до $V_2$, на изотерме с температурой $T_2'$ отводится количество тепла $Q_2$.

\problembox{B1}{Determine the amounts of heat $Q_1$ and $Q_2$, express the answer in terms of $\nu$, $R$, $T_1'$, $T_2'$, $V_1$, $V_2$.}{0.80}

\problembox{B2}{Express the cycle time $t_c$ in terms of $Q_1$, $Q_2$, $\alpha_1$, $\alpha_2$, $\Delta T_1$, $\Delta T_2$. The time taken by adiabatic processes can be neglected.}{0.40}

\problembox{B3}{Express the useful power of the cycle $P$ in terms of $T_1$, $T_2$, $\Delta T_1$, $\Delta T_2$, $\alpha_1$, $\alpha_2$.}{0.40}

In the following paragraphs we will determine the values ​​of the temperature differences at which the useful power is maximum. To do this, we will change the temperature differences $\Delta T_1$ and $\Delta T_2$. At the maximum, the corresponding change in power must be zero to the first order of smallness.

\problembox{B4}{We will consider only cycles lasting a fixed time $t_c$. Find the following relationship between the increments $d \Delta T_1$ and $d \Delta T_2$. Express your answer in terms of $T_1$, $T_2$, $\alpha_1$, $\alpha_2$, $\Delta T_1$, $\Delta T_2$.}{0.40}

\problembox{B5}{Let the cycle duration $t_c$ be fixed. Find the ratio $\Delta T_1/\Delta T_2$ at which the power $P$ is maximum.}{0.50}

\problembox{B6}{Differentiate the power by $\Delta T_1$. Show that the condition for this derivative to be zero can be written as $$ \textbackslash\{\}frac\{\textbackslash\{\}Delta T_1\}\{\textbackslash\{\}Delta T_2\} = h\textbackslash\{\}frac\{P\}\{\textbackslash\{\}Delta T_1 \textbackslash\{\}Delta T_2\}, $$где $h$ – некоторая величина, не зависящая от $\Delta T_1$ и $\Delta T_2$. Получите выражение для $h$.}{0.80}

\problembox{B7}{Using the results of the previous two paragraphs, determine the values ​​of $\Delta T_1$, $\Delta T_2$ at which the power is maximum. There is no need to prove that the values ​​found correspond exactly to the maximum. Express your answer in terms of $T_1$, $T_2$, $\alpha_1$, $\alpha_2$. You can significantly simplify the calculations if you assume that the relationship between $\Delta T_1$ and $\Delta T_2$ obtained in the previous paragraph holds true.}{1.20}

\problembox{B8}{Find the value of the maximum power $P$ and efficiency $\eta$ of the cycle at optimal values ​​$\Delta T_1$, $\Delta T_2$. Express your answer in terms of $T_1$, $T_2$, $\alpha_1$, $\alpha_2$.}{0.70}

\problembox{B9}{Typical temperature of thermal power plant heater $560 ~^\circ \mbox{C}$, refrigerator $25~^\circ \mbox{C}$. Calculate the efficiency of the Carnot cycle and the efficiency found in the previous paragraph for these temperatures. Which of these values ​​best describes the measured value $\eta = 0.36$?}{0.30}

\section*{Part C. Leaking motor pair (3 points)}

\begin{figure}[H]
    \centering
    \includegraphics[width=0.72\textwidth]{T1_EN_assets/fig1.jpg}
    \caption{Figure 1.}
\end{figure}

Due to non-ideal thermal insulation of a heat engine, it may happen that part of the heat supplied to it is not transferred to the working fluid, but directly flows from the heater to the refrigerator. To explore the role of such internal heat leaks, consider a heat engine, which is a chain of two engines. In this case, the refrigerator of the first engine is the heater of the second. Both engines can perform useful work. In addition, some of the heat is transferred directly from the engine heater to the refrigerator. Let us introduce the following notation $Q_{Hi}$ – the amount of heat that reached the working body of the i-th heat engine $Q_{Li}$ – the amount of heat that the working body of the i-th heat engine gave to its refrigerator $\eta_i$ – the efficiency of the i-th machine $Q_{Bi}$ – the amount of heat that in the i-th machine directly transferred from the heater to the refrigerator $T_H$ – the temperature of the first heater $T_{L1}$ – temperature of the first refrigerator (also the heater of the second machine) $T_L$ – temperature of the second refrigerator The first engine receives from the heater a certain amount of heat $Q_H$, which is partially transferred to the working fluid, and partially transferred directly to the refrigerators of the first and second heat engines due to leaks. We will assume that the leakage values ​​are proportional to the temperature differences of the corresponding heater and refrigerator, as well as the heat supplied to the first heater: $$Q_\{B1\} = N_1 Q_H =  K_1 (T_H - T_\{L1\}) Q_H, \textbackslash\{\}quad Q_\{B2\}=N_2 Q_H = K_2 (T_\{L1\} - T_\{L\}) Q_H$$

Введем следующие обозначения

Первый двигатель получает от нагревателя некоторое количество теплоты $Q_H$, которое частично передается рабочему телу, а частично передается непосредственно холодильникам первой и второй тепловой машин за счет утечек. Будем считать, что величины утечек пропорциональны разностям температур соответствующих нагревателя и холодильника, а также подведенному к первому нагревателю теплу:

$$Q_\{B1\} = N_1 Q_H =  K_1 (T_H - T_\{L1\}) Q_H, \textbackslash\{\}quad Q_\{B2\}=N_2 Q_H = K_2 (T_\{L1\} - T_\{L\}) Q_H$$

\problembox{C1}{Express the efficiency coefficient $\eta$ of the heat engine through the efficiency coefficients of its parts $\eta_1$, $\eta_2$ (they are defined as the ratio of useful work to the amount of heat supplied to the working fluid of the corresponding machine), as well as through the coefficients $N_1$, $N_2$.}{0.80}

\problembox{C2}{Express $\eta$ in terms of temperatures $T_H$, $T_L$, $T_{L1}$ and coefficients $K_1$, $K_2$. Think of the individual heat engine cycles as Carnot cycles.}{0.90}

\problembox{C3}{For $T_L/T_H = 0.5$, $K_1 T_H = 0.3$, $K_2 T_H = 0.2$, plot the efficiency as a function of the ratio $T_{L1}/T_H$.}{0.40}

\problembox{C4}{Find the temperature of the intermediate cooler $T_{L1}$ at which the found efficiency is maximum. For the case $K_1  = K_2$ give an expression for the maximum efficiency.}{0.90}

\vspace{1em}
\noindent\textit{Source:} \url{https://pho.rs/p/4557}

\end{document}